\chapter{Second test chapter}
\clearpage

\section{Methods}
\lipsum[7-8]

The overall processing pipeline, including a \ac{qc} step, is illustrated in Figure~\ref{fig:ch2-flow}.

\begin{figure}
  \centering
  \resizebox{\linewidth}{!}{%
  \begin{tikzpicture}[
      node distance  = 5mm and 12mm,
      box/.style     = {rectangle, draw, rounded corners = 2pt,
                        minimum width = 26mm, minimum height = 8mm,
                        text width = 24mm, align = center, font = \small},
      decision/.style= {diamond, draw, aspect = 2,
                        minimum width = 26mm,
                        text width = 18mm, align = center, font = \small},
      arr/.style     = {-Stealth, semithick},
    ]
    \node[box]      (input)  {Data input};
    \node[box,      right=of input]   (proc)   {Preprocessing};
    \node[decision, right=of proc]    (qc)     {QC passed?};
    \node[box,      right=of qc]      (output) {Analysis \& output};

    \draw[arr] (input)  -- (proc);
    \draw[arr] (proc)   -- (qc);
    \draw[arr] (qc)     -- node[above, font=\footnotesize]{yes} (output);
    \draw[arr] (qc.south) -- ++(0,-10mm) -|
               node[below, font=\footnotesize, pos=0.25]{no} (proc.south);
  \end{tikzpicture}%
  }
  \caption{Schematic overview of the data processing pipeline. Samples that
           fail the \ac{qc} check are returned to the preprocessing step.}
  \label{fig:ch2-flow}
\end{figure}

The total energy of the system is obtained by integrating the spectral density
over all frequencies~\cite{brown2021, smith2020}:
\begin{equation}
  E = \int_0^{\infty} S(f)\, df.
\end{equation}
\lipsum[9]

\section{Results and discussion}
\lipsum[10-12]

\printbibliography[heading=subbibliography, title={References}]
